\chapter{Relazioni binarie e relazioni di equivalenza}\label{cap:equivalenze}

Nel Capitolo~\ref{cap:corrispondenze} abbiamo definito le \textbf{relazioni binarie}: le corrispondenze di un insieme con
sé stesso. In questo capitolo ne studiamo le proprietà e ci concentriamo sulle più importanti, le \textbf{relazioni di
equivalenza}. Servono a raggruppare gli elementi di un insieme in \textbf{classi} di elementi che, per lo scopo che ci
interessa, consideriamo «uguali»: i numeri pari e i numeri dispari, gli orari che segnano la stessa ora sull'orologio,
gli elementi a cui una funzione assegna lo stesso valore.

% ──────────────────────────────────────────────────────────────
\section{Rappresentare una relazione binaria}\label{sec:rappresentare-relazioni}

Ricordiamo (Definizione~\ref{def:relazione-binaria}) che una relazione binaria su un insieme $s$ è una terna
$\rho = (s, s, g)$ con $g \subseteq s \times s$, e che scriviamo $x \rho y$ quando $(x, y) \in g$. Quando $s$ è finito
una relazione si può descrivere in tre modi equivalenti, ognuno utile in situazioni diverse.

\begin{esempio}{Tre modi di descrivere la stessa relazione}
    Siano $s = \{1, 2, 3\}$ e $\rho = (s, s, g)$ con $g = \{(1, 1), (1, 2), (2, 1), (2, 3)\}$.
    \begin{center}
    \begin{minipage}[c]{0.28\linewidth}\centering
        \textbf{Elenco delle coppie}\\[2mm]
        $g = \{(1,1), (1,2),$\\[0.5mm] $(2,1), (2,3)\}$
    \end{minipage}%
    \hfill
    \begin{minipage}[c]{0.28\linewidth}\centering
        \textbf{Tabella}\\[2mm]
        \renewcommand{\arraystretch}{1.25}
        \begin{tabular}{c|ccc}
            $\rho$ & $1$ & $2$ & $3$ \\ \hline
            $1$ & $\bullet$ & $\bullet$ & \\
            $2$ & $\bullet$ & & $\bullet$ \\
            $3$ & & & \\
        \end{tabular}
    \end{minipage}%
    \hfill
    \begin{minipage}[c]{0.38\linewidth}\centering
        \textbf{Grafo con frecce}\\[1mm]
        \begin{tikzpicture}
            \node[rv] (1) at (0,0) {$1$};
            \node[rv] (2) at (1.7,0) {$2$};
            \node[rv] (3) at (3.4,0) {$3$};
            \draw[rfreccia] (1) to[rcappio] (1);
            \draw[rfreccia] (1) to[bend left=28] (2);
            \draw[rfreccia] (2) to[bend left=28] (1);
            \draw[rfreccia] (2) to (3);
        \end{tikzpicture}
    \end{minipage}
    \end{center}
    Nella tabella il pallino nella \textbf{riga} di $x$ e nella \textbf{colonna} di $y$ dice che $x \rho y$. Nel grafo la stessa
    informazione è data da una \textbf{freccia} da $x$ a $y$; una freccia da un elemento a sé stesso si chiama \textbf{cappio}.
    Qui $1 \rho 1$, $1 \rho 2$, $2 \rho 1$ e $2 \rho 3$, mentre ad esempio $3 \rho 2$ è falso.
\end{esempio}

% ──────────────────────────────────────────────────────────────
\section{Proprietà delle relazioni binarie}\label{sec:proprieta-relazioni}

\dfn[proprieta-relazioni]{Proprietà delle relazioni binarie}{
    Sia $\rho = (s, s, g)$ una relazione binaria. $\rho$ si dice
    \begin{itemize}
        \item \textbf{riflessiva} se $\forall x \in s\,(x \rho x)$;
        \item \textbf{antiriflessiva} se $\forall x \in s\,\neg(x \rho x)$;
        \item \textbf{simmetrica} se $\forall x, y \in s\,(x \rho y \rightarrow y \rho x)$;
        \item \textbf{asimmetrica} se $\forall x, y \in s\,\big(x \rho y \rightarrow \neg(y \rho x)\big)$;
        \item \textbf{antisimmetrica} se $\forall x, y \in s\,\big((x \rho y \wedge y \rho x) \rightarrow x = y\big)$;
        \item \textbf{transitiva} se $\forall x, y, z \in s\,\big((x \rho y \wedge y \rho z) \rightarrow x \rho z\big)$.
    \end{itemize}
}

\begin{concetto}
Sul grafo con frecce le proprietà si riconoscono a colpo d'occhio:
\begin{center}
\renewcommand{\arraystretch}{1.35}
\begin{tabular}{>{\bfseries}l p{0.66\linewidth}}
    \rowcolor{primary!12} Proprietà & \textbf{Come si vede sul grafo} \\
    riflessiva & c'è un cappio su \textbf{ogni} elemento \\
    antiriflessiva & non c'è \textbf{nessun} cappio \\
    simmetrica & ogni freccia ha la sua \textbf{freccia di ritorno} \\
    asimmetrica & \textbf{nessuna} freccia ha la freccia di ritorno (in particolare non ci sono cappi) \\
    antisimmetrica & tra due elementi \textbf{distinti} non ci sono mai frecce in entrambi i versi (i cappi sono ammessi) \\
    transitiva & ogni volta che si va da $x$ a $y$ e da $y$ a $z$, c'è anche la \textbf{scorciatoia} da $x$ a $z$ \\
\end{tabular}
\end{center}
\end{concetto}

\dfn[tipi-relazioni]{Relazioni di equivalenza e relazioni d'ordine}{
    Sia $\rho$ una relazione binaria su $s$. $\rho$ si dice
    \begin{itemize}
        \item \textbf{relazione di equivalenza} se è riflessiva, simmetrica e transitiva;
        \item \textbf{relazione d'ordine largo} (o semplicemente \textbf{d'ordine}) se è riflessiva, antisimmetrica e transitiva;
        \item \textbf{relazione d'ordine stretto} se è antiriflessiva e transitiva.
    \end{itemize}
    Indichiamo con $\mathrm{Eq}(s)$ l'insieme delle relazioni di equivalenza su $s$.
    \begin{center}
    \renewcommand{\arraystretch}{1.25}
    \begin{tabular}{l|ccccc}
        & riflessiva & antiriflessiva & simmetrica & antisimmetrica & transitiva \\ \hline
        equivalenza & \checkmark & & \checkmark & & \checkmark \\
        ordine largo & \checkmark & & & \checkmark & \checkmark \\
        ordine stretto & & \checkmark & & & \checkmark \\
    \end{tabular}
    \end{center}
}

\begin{esempio}{Relazioni notevoli}
    Sia $s$ un insieme.
    \begin{itemize}
        \item La relazione \textbf{identica} $(s, s, \mathrm{diag}(s))$, cioè l'uguaglianza ($x \rho y \leftrightarrow x = y$), è sia
              una relazione di equivalenza sia una relazione d'ordine largo. Non è d'ordine stretto, a meno che $s = \varnothing$.
        \item La relazione \textbf{totale} $(s, s, s \times s)$ è una relazione di equivalenza. È antiriflessiva solo se
              $s = \varnothing$ e antisimmetrica solo se $s$ ha al più un elemento.
        \item L'inclusione $\subseteq$ è una relazione d'ordine largo su $\mcP(s)$, l'inclusione stretta $\subset$ una relazione
              d'ordine stretto.
        \item $\leq$ è una relazione d'ordine largo su $\mbZ$ e $<$ una relazione d'ordine stretto.
        \item La relazione di divisibilità su $\mbN$, $m \mid n \leftrightarrow \exists k \in \mbN\,(n = km)$, è una relazione
              d'ordine largo.
        \item \textbf{Attenzione}: la divisibilità su $\mbZ$, $m \mid n \leftrightarrow \exists k \in \mbZ\,(n = km)$,
              \textbf{non} è una relazione d'ordine, perché non è antisimmetrica: $1 \mid -1$ e $-1 \mid 1$, ma $1 \neq -1$.
    \end{itemize}
\end{esempio}

Per dimostrare che una proprietà \textbf{non} vale bisogna negarne la definizione. Le negazioni hanno sempre la forma
«esiste un controesempio».

\begin{concetto}
\begin{Proposizione}{Negazione delle proprietà}{negazioni-relazioni}
    \begin{prerequisiti}
        \dip{Definizione}{def:proprieta-relazioni}
        \dip{Teorema}{teo:negazione-quantificatori}
        \dip{Proposizione}{prop:negimp}
        \dip{Proposizione}{prop:demorgan}
    \end{prerequisiti}
    Sia $\rho$ una relazione binaria su $s$. Allora:
    \begin{enumerate}
        \item $\rho$ \textbf{non} è riflessiva $\leftrightarrow \exists x \in s\,\neg(x \rho x)$;
        \item $\rho$ \textbf{non} è simmetrica $\leftrightarrow \exists x, y \in s\,\big(x \rho y \wedge \neg(y \rho x)\big)$;
        \item $\rho$ \textbf{non} è antisimmetrica $\leftrightarrow \exists x, y \in s\,(x \rho y \wedge y \rho x \wedge x \neq y)$;
        \item $\rho$ \textbf{non} è transitiva $\leftrightarrow \exists x, y, z \in s\,\big(x \rho y \wedge y \rho z \wedge \neg(x \rho z)\big)$.
    \end{enumerate}
\end{Proposizione}
\begin{dimostrazione}
    In ogni punto si nega la definizione: il $\forall$ diventa $\exists$ e si nega ciò che segue
    (Teorema~\ref{teo:negazione-quantificatori}); poi si usa $\neg(p \rightarrow q) \leftrightarrow (p \wedge \neg q)$
    (Proposizione~\ref{prop:negimp}).
    \begin{caso}{Punto 1}
        $\neg\forall x\,(x \rho x) \leftrightarrow \exists x\,\neg(x \rho x)$.
    \end{caso}
    \begin{caso}{Punti 2, 3, 4}
        Per il punto 2: $\neg\forall x, y\,(x \rho y \rightarrow y \rho x) \leftrightarrow \exists x, y\,\neg(x \rho y \rightarrow y \rho x)
        \leftrightarrow \exists x, y\,(x \rho y \wedge \neg(y \rho x))$. Per il punto 3 si nega l'implicazione
        $(x \rho y \wedge y \rho x) \rightarrow x = y$ e si ottiene $x \rho y \wedge y \rho x \wedge \neg(x = y)$; per il punto 4
        si nega $(x \rho y \wedge y \rho z) \rightarrow x \rho z$ e si ottiene $x \rho y \wedge y \rho z \wedge \neg(x \rho z)$.
    \end{caso}
\end{dimostrazione}
\end{concetto}

\begin{spiegazione}
    Per mostrare che una proprietà \textbf{vale} bisogna controllare \textbf{tutti} gli elementi; per mostrare che \textbf{non vale}
    basta \textbf{un solo controesempio}. Sul grafo i quattro controesempi sono questi (in rosso tratteggiato ciò che manca o che
    non dovrebbe esserci):
    \begin{center}
    \begin{tikzpicture}
        \begin{scope}
            \node[rv] (x) at (0,0) {$x$};
            \draw[rfreccia, cherryred, dashed] (x) to[rcappio] (x);
            \node[font=\small, align=center] at (0,-0.95) {non riflessiva:\\manca un cappio};
        \end{scope}
        \begin{scope}[xshift=3.2cm]
            \node[rv] (x) at (0,0) {$x$}; \node[rv] (y) at (1.4,0) {$y$};
            \draw[rfreccia] (x) to[bend left=25] (y);
            \draw[rfreccia, cherryred, dashed] (y) to[bend left=25] (x);
            \node[font=\small, align=center] at (0.7,-0.95) {non simmetrica:\\manca il ritorno};
        \end{scope}
        \begin{scope}[xshift=7.4cm]
            \node[rv] (x) at (0,0) {$x$}; \node[rv] (y) at (1.4,0) {$y$};
            \draw[rfreccia, cherryred] (x) to[bend left=25] (y);
            \draw[rfreccia, cherryred] (y) to[bend left=25] (x);
            \node[font=\small, align=center] at (0.7,-0.95) {non antisimmetrica:\\andata e ritorno, $x \neq y$};
        \end{scope}
        \begin{scope}[xshift=11.6cm]
            \node[rv] (x) at (0,0) {$x$}; \node[rv] (y) at (0.8,0.95) {$y$}; \node[rv] (z) at (1.6,0) {$z$};
            \draw[rfreccia] (x) to (y); \draw[rfreccia] (y) to (z);
            \draw[rfreccia, cherryred, dashed] (x) to (z);
            \node[font=\small, align=center] at (0.8,-0.95) {non transitiva:\\manca la scorciatoia};
        \end{scope}
    \end{tikzpicture}
    \end{center}
\end{spiegazione}

\warning{Negare non vuol dire passare alla proprietà «anti»}{
    «Non riflessiva» \textbf{non} significa «antiriflessiva», e «non simmetrica» non significa «antisimmetrica».
    Su $s = \{1, 2, 3\}$:
    \begin{itemize}
        \item $g = \{(1, 1)\}$ non è riflessiva ($2$ non ha il cappio) e non è antiriflessiva ($1$ ce l'ha);
        \item $g = \{(1, 2), (2, 1), (2, 3)\}$ non è simmetrica ($2 \rho 3$ ma non $3 \rho 2$) e non è antisimmetrica
              ($1 \rho 2$ e $2 \rho 1$ con $1 \neq 2$);
        \item l'uguaglianza è \textbf{sia} simmetrica \textbf{sia} antisimmetrica.
    \end{itemize}
}

Controllare la transitività su un insieme finito richiede di guardare tutte le terne di elementi. Molte terne però si
controllano da sole.

\begin{concetto}
\begin{Proposizione}{Quali terne guardare per la transitività}{transitivita-distinti}
    \begin{prerequisiti}
        \dip{Definizione}{def:proprieta-relazioni}
        \dip{Proposizione}{prop:terzoescluso}
    \end{prerequisiti}
    Sia $\rho$ una relazione binaria su $s$. $\rho$ è transitiva se e solo se valgono entrambe le condizioni:
    \begin{enumerate}
        \item[(a)] $(x \rho y \wedge y \rho z) \rightarrow x \rho z$ per ogni terna $x, y, z$ di elementi \textbf{a due a due distinti};
        \item[(b)] $(x \rho y \wedge y \rho x) \rightarrow x \rho x$ per ogni $x, y \in s$.
    \end{enumerate}
    In particolare, se $\rho$ è \textbf{riflessiva} oppure \textbf{antisimmetrica}, la (b) vale sempre, e per la transitività
    \textbf{basta guardare le terne di elementi a due a due distinti}.
\end{Proposizione}
\begin{dimostrazione}
    \begin{caso}{$(\rightarrow)$ Se $\rho$ è transitiva, valgono (a) e (b)}
        (a) è un caso particolare della transitività; (b) è la transitività applicata alla terna $x, y, x$.
    \end{caso}
    \begin{caso}{$(\leftarrow)$ Se valgono (a) e (b), $\rho$ è transitiva}
        Siano $x \rho y$ e $y \rho z$: dobbiamo mostrare $x \rho z$.
        Se $x, y, z$ sono a due a due distinti, basta (a). Se $x = y$, la tesi $x \rho z$ è l'ipotesi $y \rho z$; se $y = z$,
        la tesi è l'ipotesi $x \rho y$. Resta il caso $x = z$ (con $x \neq y$): le ipotesi sono $x \rho y$ e $y \rho x$, e
        la tesi $x \rho x$ è proprio la (b).
    \end{caso}
    \begin{caso}{Se $\rho$ è riflessiva o antisimmetrica, la (b) vale sempre}
        Se $\rho$ è riflessiva, $x \rho x$ vale comunque. Se $\rho$ è antisimmetrica, da $x \rho y$ e $y \rho x$ segue $x = y$,
        e allora $x \rho x$ è l'ipotesi $x \rho y$.
    \end{caso}
\end{dimostrazione}
\end{concetto}

\begin{spiegazione}
    In una terna con due elementi uguali la transitività è quasi sempre automatica: se $x = y$ o $y = z$ la «scorciatoia»
    coincide con una delle due frecce che abbiamo già. L'unico caso insidioso è il \textbf{viaggio di andata e ritorno}
    $x \to y \to x$, che obbliga ad avere il cappio su $x$. Ecco perché la condizione (b) non si può togliere:
    \begin{center}
    \begin{tikzpicture}
        \node[rv] (1) at (0,0) {$1$}; \node[rv] (2) at (1.8,0) {$2$};
        \draw[rfreccia] (1) to[bend left=28] (2);
        \draw[rfreccia] (2) to[bend left=28] (1);
        \draw[rfreccia, cherryred, dashed] (1) to[rcappio] (1);
        \draw[rfreccia, cherryred, dashed] (2) to[rcappio] (2);
        \node[font=\small, align=left, anchor=west] at (2.8,0.1)
            {$s = \{1, 2\}$, $g = \{(1, 2), (2, 1)\}$: non ci sono terne di elementi distinti,\\
             quindi (a) è vera; ma $1 \rho 2$ e $2 \rho 1$ senza $1 \rho 1$: $\rho$ \textbf{non} è transitiva.};
    \end{tikzpicture}
    \end{center}
    Per le relazioni di equivalenza e d'ordine largo (riflessive) e per quelle antisimmetriche questo problema non c'è.
\end{spiegazione}

La definizione di ordine stretto chiede solo antiriflessività e transitività. Ci si aspetterebbe di dover chiedere anche
che non ci siano mai «andata e ritorno», cioè l'asimmetria: la proposizione seguente mostra che viene gratis.

\begin{concetto}
\begin{Proposizione}{Per un ordine stretto bastano antiriflessività e transitività}{stretto-asimmetrico}
    \begin{prerequisiti}
        \dip{Definizione}{def:proprieta-relazioni}
        \dip{Definizione}{def:tipi-relazioni}
    \end{prerequisiti}
    Sia $\rho$ una relazione binaria transitiva su $s$. Allora $\rho$ è asimmetrica se e solo se è antiriflessiva.
    In particolare ogni relazione d'ordine stretto è asimmetrica e antisimmetrica.
\end{Proposizione}
\begin{dimostrazione}
    \begin{caso}{$(\rightarrow)$ Asimmetrica $\Rightarrow$ antiriflessiva (anche senza transitività)}
        Se fosse $x \rho x$ per qualche $x$, l'asimmetria con $y = x$ darebbe $\neg(x \rho x)$: assurdo.
    \end{caso}
    \begin{caso}{$(\leftarrow)$ Antiriflessiva e transitiva $\Rightarrow$ asimmetrica}
        Per assurdo siano $x \rho y$ e $y \rho x$. Per transitività $x \rho x$, contro l'antiriflessività.
    \end{caso}
    \begin{caso}{Ogni ordine stretto è antisimmetrico}
        È asimmetrico per quanto appena visto, quindi $x \rho y \wedge y \rho x$ è sempre falsa e l'implicazione che definisce
        l'antisimmetria è vera (antecedente falso).
    \end{caso}
\end{dimostrazione}
\end{concetto}

\begin{spiegazione}
    Un'«andata e ritorno» $x \to y \to x$, per la transitività, produrrebbe la scorciatoia $x \to x$: un cappio. Ma un ordine
    stretto non ha cappi, quindi non può avere andate e ritorni. Pensiamo a $<$ su $\mbZ$: da $3 < 5$ e $5 < 3$ seguirebbe
    $3 < 3$, che è falso; quindi $3 < 5$ e $5 < 3$ non possono essere vere insieme.
\end{spiegazione}

% ──────────────────────────────────────────────────────────────
\section{Relazione duale e diagonale}\label{sec:duale-diagonale}

\dfn[relazione-duale]{Relazione duale}{
    Sia $\rho = (s, s, g)$ una relazione binaria. La relazione $\bar{\rho} = (s, s, \bar{g})$ con
    \[ \bar{g} := \{(y, x) \mid (x, y) \in g\}, \qquad\text{cioè}\qquad \forall x, y \in s\,(x \bar{\rho} y \leftrightarrow y \rho x), \]
    si dice \textbf{relazione duale} (o \textbf{opposta}) di $\rho$. Sul grafo si ottiene \textbf{invertendo il verso di tutte le frecce}.
}

\begin{esempio}{Relazioni duali}
    La duale di $\leq$ su $\mbZ$ è $\geq$; la duale di $\subseteq$ è $\supseteq$; la duale di «essere padre di» è «essere figlio di»
    (ristretta ai maschi). La duale di una relazione simmetrica è la relazione stessa.
\end{esempio}

Ricordiamo che $\mathrm{diag}(s) = \{(x, x) \mid x \in s\}$ (Definizione~\ref{def:relazione-binaria}): nel grafo corrisponde
ai cappi, nella tabella alla diagonale. Le proprietà della Definizione~\ref{def:proprieta-relazioni} si possono scrivere come
semplici relazioni tra insiemi.

\begin{concetto}
\begin{Proposizione}{Le proprietà guardando il grafico}{proprieta-grafico}
    \begin{prerequisiti}
        \dip{Definizione}{def:proprieta-relazioni}
        \dip{Definizione}{def:relazione-duale}
        \dip{Definizione}{def:relazione-binaria}
        \dip{Definizione}{def:prodotto-corrispondenze}
        \dip{Proposizione}{prop:doppia-inclusione}
    \end{prerequisiti}
    Sia $\rho = (s, s, g)$ una relazione binaria e sia $\bar{g}$ il grafico della duale. Allora
    \begin{enumerate}
        \item $\rho$ è riflessiva $\leftrightarrow \mathrm{diag}(s) \subseteq g$;
        \item $\rho$ è antiriflessiva $\leftrightarrow \mathrm{diag}(s) \cap g = \varnothing$;
        \item $\rho$ è simmetrica $\leftrightarrow g = \bar{g}$, cioè $\rho = \bar{\rho}$;
        \item $\rho$ è antisimmetrica $\leftrightarrow g \cap \bar{g} \subseteq \mathrm{diag}(s)$;
        \item $\rho$ è asimmetrica $\leftrightarrow g \cap \bar{g} = \varnothing$;
        \item $\rho$ è transitiva $\leftrightarrow$ il grafico di $\rho\rho$ è contenuto in $g$.
    \end{enumerate}
\end{Proposizione}
\begin{dimostrazione}
    Basta tradurre le definizioni, ricordando che $(x, y) \in \bar{g} \leftrightarrow (y, x) \in g$.
    \begin{caso}{Punti 1 e 2}
        $\mathrm{diag}(s) \subseteq g$ significa che ogni $(x, x)$ sta in $g$, cioè $x \rho x$ per ogni $x$. $\mathrm{diag}(s) \cap g = \varnothing$
        significa che nessun $(x, x)$ sta in $g$, cioè $\neg(x \rho x)$ per ogni $x$.
    \end{caso}
    \begin{caso}{Punto 3}
        $(\rightarrow)$ Se $(x, y) \in g$, per simmetria $(y, x) \in g$, cioè $(x, y) \in \bar{g}$: quindi $g \subseteq \bar{g}$. Se
        $(x, y) \in \bar{g}$, allora $(y, x) \in g$ e per simmetria $(x, y) \in g$: quindi $\bar{g} \subseteq g$, e per la doppia
        inclusione $g = \bar{g}$. $(\leftarrow)$ Se $x \rho y$, allora $(x, y) \in g = \bar{g}$, cioè $(y, x) \in g$: $y \rho x$.
    \end{caso}
    \begin{caso}{Punti 4 e 5}
        $(x, y) \in g \cap \bar{g}$ significa $x \rho y \wedge y \rho x$. L'antisimmetria dice che questo accade solo se $x = y$, cioè solo
        se $(x, y) \in \mathrm{diag}(s)$; l'asimmetria dice che non accade mai, cioè $g \cap \bar{g} = \varnothing$.
    \end{caso}
    \begin{caso}{Punto 6}
        Per la Definizione~\ref{def:prodotto-corrispondenze}, $(x, z)$ sta nel grafico di $\rho\rho$ se e solo se esiste $y$ con
        $x \rho y$ e $y \rho z$. Dire che tutte queste coppie stanno in $g$ è esattamente la transitività.
    \end{caso}
\end{dimostrazione}
\end{concetto}

\begin{spiegazione}
    Conviene pensare alla \textbf{tabella}. La diagonale $\mathrm{diag}(s)$ sono le caselle in cui riga e colonna coincidono,
    e la duale si ottiene \textbf{ribaltando la tabella attorno alla diagonale}: la casella (riga $x$, colonna $y$) va in
    (riga $y$, colonna $x$).
    \begin{center}
    \begin{tikzpicture}[cella/.style={minimum size=6mm, draw=black!25, font=\small}]
        \foreach \i in {1,2,3} {
            \node[font=\small] at (\i*0.6, 0.6) {$\i$};
            \node[font=\small] at (0, -\i*0.6+0.6) {$\i$};
        }
        \foreach \r/\c in {1/1,2/2,3/3} \node[cella, fill=naplesyellow!40] at (\c*0.6, -\r*0.6+0.6) {};
        \foreach \r/\c in {1/2,2/1,2/3} \node[cella] at (\c*0.6, -\r*0.6+0.6) {$\bullet$};
        \foreach \r/\c in {1/1} \node[cella, fill=naplesyellow!40] at (\c*0.6, -\r*0.6+0.6) {$\bullet$};
        \node[font=\small] at (1.2,-1.75) {$g$};
        \draw[-{Stealth}, line width=0.8pt, primary] (2.6,-0.6) -- node[above, font=\small] {ribalto} (3.8,-0.6);
        \begin{scope}[xshift=4.2cm]
            \foreach \i in {1,2,3} {
                \node[font=\small] at (\i*0.6, 0.6) {$\i$};
                \node[font=\small] at (0, -\i*0.6+0.6) {$\i$};
            }
            \foreach \r/\c in {1/1,2/2,3/3} \node[cella, fill=naplesyellow!40] at (\c*0.6, -\r*0.6+0.6) {};
            \foreach \r/\c in {2/1,1/2,3/2} \node[cella] at (\c*0.6, -\r*0.6+0.6) {$\bullet$};
            \foreach \r/\c in {1/1} \node[cella, fill=naplesyellow!40] at (\c*0.6, -\r*0.6+0.6) {$\bullet$};
            \node[font=\small] at (1.2,-1.75) {$\bar{g}$};
        \end{scope}
        \node[font=\small, align=left, anchor=west] at (7.0,-0.55)
            {È la relazione dell'Esempio iniziale:\\
             $g = \{(1,1), (1,2), (2,1), (2,3)\}$.\\
             La diagonale (in giallo) resta ferma.};
    \end{tikzpicture}
    \end{center}
    Allora: \textbf{riflessiva} = diagonale tutta piena; \textbf{antiriflessiva} = diagonale tutta vuota; \textbf{simmetrica} =
    la tabella non cambia ribaltandola; \textbf{antisimmetrica} = fuori dalla diagonale, una casella e la sua simmetrica non sono
    mai piene insieme. Nell'esempio $(1, 2)$ e $(2, 1)$ sono piene insieme: la relazione non è antisimmetrica; e $(2, 3)$ è
    piena mentre $(3, 2)$ no: non è nemmeno simmetrica.
\end{spiegazione}

\begin{concetto}
\begin{Corollario}{Relazioni simmetriche e antisimmetriche}{simm-antisimm}
    \begin{prerequisiti}
        \dip{Proposizione}{prop:proprieta-grafico}
    \end{prerequisiti}
    Una relazione binaria $\rho = (s, s, g)$ è sia simmetrica sia antisimmetrica se e solo se $g \subseteq \mathrm{diag}(s)$.
\end{Corollario}
\begin{dimostrazione}
    \begin{caso}{$(\rightarrow)$}
        Per i punti 3 e 4 della Proposizione~\ref{prop:proprieta-grafico}, $g = \bar{g}$, quindi $g = g \cap \bar{g} \subseteq \mathrm{diag}(s)$.
    \end{caso}
    \begin{caso}{$(\leftarrow)$}
        Se $g \subseteq \mathrm{diag}(s)$, ogni coppia di $g$ è del tipo $(x, x)$, che coincide con la sua ribaltata: quindi
        $\bar{g} = g$ e $\rho$ è simmetrica. Inoltre $g \cap \bar{g} \subseteq g \subseteq \mathrm{diag}(s)$: $\rho$ è antisimmetrica.
    \end{caso}
\end{dimostrazione}
\end{concetto}

\begin{spiegazione}
    Simmetrica vuol dire «ogni freccia ha il ritorno», antisimmetrica «nessuna freccia tra elementi distinti ha il ritorno».
    Le due cose insieme sono possibili solo se \textbf{non ci sono frecce tra elementi distinti}: restano al più i cappi.
    Ad esempio su $\{1, 2, 3\}$ la relazione con grafico $\{(1,1), (3,3)\}$ è simmetrica e antisimmetrica.
\end{spiegazione}

\begin{concetto}
\begin{Proposizione}{La duale conserva le proprietà}{duale-proprieta}
    \begin{prerequisiti}
        \dip{Definizione}{def:relazione-duale}
        \dip{Definizione}{def:proprieta-relazioni}
        \dip{Definizione}{def:tipi-relazioni}
        \dip{Corollario}{cor:terna}
    \end{prerequisiti}
    Sia $\rho$ una relazione binaria su $s$. Allora $\bar{\bar{\rho}} = \rho$ e $\bar{\rho}$ ha le stesse proprietà di $\rho$ tra quelle
    della Definizione~\ref{def:proprieta-relazioni}. In particolare la duale di una relazione di equivalenza, d'ordine largo o
    d'ordine stretto è una relazione dello stesso tipo.
\end{Proposizione}
\begin{dimostrazione}
    \begin{caso}{$\bar{\bar{\rho}} = \rho$}
        $x \bar{\bar{\rho}} y \leftrightarrow y \bar{\rho} x \leftrightarrow x \rho y$: i grafici coincidono, quindi (Corollario~\ref{cor:terna}) le terne coincidono.
    \end{caso}
    \begin{caso}{Le proprietà si conservano}
        Ogni proprietà si ottiene da quella di $\rho$ \textbf{scambiando i nomi} delle variabili. Ad esempio, per la transitività:
        se $x \bar{\rho} y$ e $y \bar{\rho} z$, allora $z \rho y$ e $y \rho x$, quindi $z \rho x$ per la transitività di $\rho$, cioè
        $x \bar{\rho} z$. Per l'antisimmetria: $x \bar{\rho} y \wedge y \bar{\rho} x$ è $y \rho x \wedge x \rho y$, che implica $x = y$.
        Riflessività e antiriflessività riguardano $x \bar{\rho} x$, che è $x \rho x$; simmetria e asimmetria si trattano allo stesso modo.
        Viceversa, se $\bar{\rho}$ ha una proprietà, ce l'ha anche $\bar{\bar{\rho}} = \rho$.
    \end{caso}
\end{dimostrazione}
\end{concetto}

\begin{spiegazione}
    Invertire tutte le frecce non cambia la «forma» del grafo: i cappi restano cappi, le andate e ritorno restano andate e
    ritorno, e un percorso $x \to y \to z$ con la sua scorciatoia $x \to z$ diventa il percorso $z \to y \to x$ con la scorciatoia
    $z \to x$. Ad esempio $\geq$ è un ordine largo come $\leq$, e $>$ un ordine stretto come $<$.
\end{spiegazione}

% ──────────────────────────────────────────────────────────────
\section{Relazioni di equivalenza}\label{sec:equivalenze}

Una relazione di equivalenza si comporta come una «uguaglianza parziale»: ogni elemento è equivalente a sé stesso, se $x$ è
equivalente a $y$ anche $y$ lo è a $x$, e due elementi equivalenti a un terzo sono equivalenti tra loro.

\begin{esempio}{Relazioni di equivalenza}
    \begin{itemize}
        \item L'uguaglianza e la relazione totale su un qualunque insieme.
        \item Sull'insieme delle persone, «avere la stessa età» o «essere nati nello stesso mese».
        \item In un insieme di strutture algebriche, «essere isomorfe»: ogni struttura è isomorfa a sé stessa tramite l'identità,
              l'inversa di un isomorfismo è un isomorfismo (Teorema~\ref{teo:inverso-isomorfismo}) e la composta di due
              isomorfismi è biettiva (Corollario~\ref{cor:comp-biettive}) e rispetta l'operazione.
    \end{itemize}
\end{esempio}

\dfn[congruenza-modulo]{Congruenza modulo $m$}{
    Sia $m \in \mbZ$. Definiamo la relazione binaria $\equiv_m$ su $\mbZ$ ponendo
    \[ \forall a, b \in \mbZ\,\big(a \equiv_m b \leftrightarrow \exists k \in \mbZ\,(b - a = km)\big). \]
    Se $a \equiv_m b$ diciamo che \textbf{$a$ è congruo a $b$ modulo $m$}: $a$ e $b$ differiscono per un multiplo di $m$.
}

\begin{concetto}
\begin{Proposizione}{La congruenza modulo $m$ è di equivalenza}{congruenza-equivalenza}
    \begin{prerequisiti}
        \dip{Definizione}{def:congruenza-modulo}
        \dip{Definizione}{def:tipi-relazioni}
    \end{prerequisiti}
    Per ogni $m \in \mbZ$, $\equiv_m$ è una relazione di equivalenza su $\mbZ$.
\end{Proposizione}
\begin{dimostrazione}
    \begin{caso}{Riflessiva}
        Per ogni $a \in \mbZ$, $a - a = 0 = 0 \cdot m$, quindi $a \equiv_m a$ (con $k = 0$).
    \end{caso}
    \begin{caso}{Simmetrica}
        Se $a \equiv_m b$, esiste $k$ con $b - a = km$. Allora $a - b = (-k)m$, quindi $b \equiv_m a$ (con $-k$).
    \end{caso}
    \begin{caso}{Transitiva}
        Se $a \equiv_m b$ e $b \equiv_m c$, esistono $h, k$ con $b - a = hm$ e $c - b = km$. Sommando,
        $c - a = (c - b) + (b - a) = km + hm = (h + k)m$, quindi $a \equiv_m c$ (con $h + k$).
    \end{caso}
\end{dimostrazione}
\end{concetto}

\begin{spiegazione}
    Pensiamo ai numeri interi come a una strada e a $\equiv_m$ come a \textbf{salti di lunghezza $m$}: $a \equiv_m b$ quando da $a$ si
    arriva a $b$ con un certo numero $k$ di salti (in avanti se $k > 0$, all'indietro se $k < 0$). Per $m = 3$, colorando ogni
    numero secondo i salti che lo raggiungono partendo da $0$, $1$ o $2$:
    \begin{center}
    \begin{tikzpicture}[x=0.92cm]
        \draw[-{Stealth}, black!50] (-4.6,0) -- (8.6,0);
        \foreach \n in {-4,...,8} {
            \pgfmathtruncatemacro{\r}{mod(\n+300,3)}
            \ifcase\r \def\col{classeA}\or \def\col{classeB}\or \def\col{classeC}\fi
            \node[circle, fill=\col, draw=black!30, inner sep=0pt, minimum size=6mm, font=\small] at (\n,0) {$\n$};
        }
        \draw[-{Stealth}, line width=0.7pt, orange!80!black] (1,0.35) to[bend left=55] node[above, font=\scriptsize] {$+3$} (4,0.35);
        \draw[-{Stealth}, line width=0.7pt, orange!80!black] (4,0.35) to[bend left=55] node[above, font=\scriptsize] {$+3$} (7,0.35);
        \draw[-{Stealth}, line width=0.7pt, orange!80!black] (1,-0.35) to[bend left=55] node[below, font=\scriptsize] {$-3$} (-2,-0.35);
    \end{tikzpicture}
    \end{center}
    \begin{passi}
        \item \textbf{Riflessiva}: da $a$ si arriva ad $a$ con zero salti.
        \item \textbf{Simmetrica}: se da $1$ si arriva a $7$ con due salti avanti, da $7$ si torna a $1$ con due salti indietro.
        \item \textbf{Transitiva}: da $1$ a $4$ con un salto e da $4$ a $7$ con un altro: da $1$ a $7$ con $1 + 1 = 2$ salti.
    \end{passi}
    I numeri dello stesso colore sono tutti congrui tra loro modulo $3$.
\end{spiegazione}

\begin{esempio}{Casi particolari}
    $\equiv_0$ è l'uguaglianza ($b - a = k \cdot 0$ vuol dire $a = b$); $\equiv_1$ è la relazione totale (ogni differenza è un
    multiplo di $1$); due interi sono congrui modulo $2$ se e solo se sono entrambi pari o entrambi dispari. Infine
    $\equiv_m$ e $\equiv_{-m}$ sono la stessa relazione, perché $km = (-k)(-m)$.
\end{esempio}

\dfn[nucleo]{Nucleo di equivalenza di una funzione}{
    Sia $f \colon a \to b$ una funzione. La relazione binaria $\sim_f$ su $a$ definita da
    \[ \forall x, y \in a\,\big(x \sim_f y \leftrightarrow f(x) = f(y)\big) \]
    si chiama \textbf{nucleo di equivalenza} di $f$ e si indica anche con $\mathrm{Ker}\, f$.
}

\begin{concetto}
\begin{Proposizione}{Proprietà del nucleo di equivalenza}{nucleo-equivalenza}
    \begin{prerequisiti}
        \dip{Definizione}{def:nucleo}
        \dip{Definizione}{def:tipi-relazioni}
        \dip{Proposizione}{prop:forme-iniettivita}
        \dip{Definizione}{def:funzioni-notevoli}
    \end{prerequisiti}
    Sia $f \colon a \to b$ una funzione. Allora:
    \begin{enumerate}
        \item $\sim_f$ è una relazione di equivalenza su $a$;
        \item $\sim_f$ è la relazione identica se e solo se $f$ è iniettiva;
        \item se $a \neq \varnothing$, $\sim_f$ è la relazione totale se e solo se $f$ è costante.
    \end{enumerate}
\end{Proposizione}
\begin{dimostrazione}
    \begin{caso}{Punto 1}
        Le tre proprietà vengono da quelle dell'uguaglianza: $f(x) = f(x)$; se $f(x) = f(y)$ allora $f(y) = f(x)$; se
        $f(x) = f(y)$ e $f(y) = f(z)$ allora $f(x) = f(z)$.
    \end{caso}
    \begin{caso}{Punto 2}
        Sempre $x = y \rightarrow f(x) = f(y)$. Quindi $\sim_f$ è l'uguaglianza se e solo se vale anche
        $f(x) = f(y) \rightarrow x = y$ per ogni $x, y$, cioè (Proposizione~\ref{prop:forme-iniettivita}) se e solo se $f$ è iniettiva.
    \end{caso}
    \begin{caso}{Punto 3}
        $(\leftarrow)$ Se $f$ è costante di valore $c$, $f(x) = c = f(y)$ per ogni $x, y$: tutti sono in relazione.
        $(\rightarrow)$ Prendo $x_0 \in a$ e pongo $c = f(x_0)$. Per ogni $x \in a$ si ha $x \sim_f x_0$, cioè $f(x) = c$: $f$ è la
        funzione costante di valore $c$.
    \end{caso}
\end{dimostrazione}
\end{concetto}

\begin{spiegazione}
    $\sim_f$ mette insieme gli elementi che \textbf{$f$ non sa distinguere}, cioè che hanno la stessa immagine. Prendiamo ad
    esempio la funzione che manda ogni intero nel suo quadrato, $f(n) = n^2$:
    \begin{center}
    \begin{tikzpicture}
        \foreach \n/\y in {0/2.4, 1/1.6, -1/1.15, 2/0.35, -2/-0.1, 3/-0.9, -3/-1.35}
            \node[elem] (a\n) at (0,\y) {$\n$};
        \foreach \v/\y in {0/2.4, 1/1.375, 4/0.125, 9/-1.125}
            \node[elem] (b\v) at (3.4,\y) {$\v$};
        \blocco{classeA}{(a0.center) -- ++(0.001,0)}
        \blocco{classeB}{(a1.center) -- (a-1.center)}
        \blocco{classeC}{(a2.center) -- (a-2.center)}
        \blocco{classeD}{(a3.center) -- (a-3.center)}
        \foreach \n/\v in {0/0, 1/1, -1/1, 2/4, -2/4, 3/9, -3/9} \draw[fun] (a\n) -- (b\v);
        \node[font=\small] at (0,3.05) {$\mbZ$ (una parte)};
        \node[font=\small] at (3.4,3.05) {$\mbZ$};
        \node[font=\small, primary!80!black] at (1.7,2.75) {$f$};
        \node[font=\small, align=left, anchor=west] at (4.4,0.6)
            {$-2 \sim_f 2$ perché $f(-2) = 4 = f(2)$;\\
             $1 \not\sim_f 2$ perché $f(1) = 1 \neq 4 = f(2)$.\\[1mm]
             Ogni macchia colorata è un gruppo di elementi\\
             con la stessa immagine.};
    \end{tikzpicture}
    \end{center}
    Il punto 2 dice che se $f$ è iniettiva non ci sono gruppi con più di un elemento; il punto 3 che se $f$ è costante c'è un
    solo gruppo, con dentro tutto $a$.
\end{spiegazione}

% ──────────────────────────────────────────────────────────────
\section{Classi di equivalenza e insieme quoziente}\label{sec:classi}

\dfn[classe-equivalenza]{Classe di equivalenza, insieme quoziente}{
    Sia $\sim$ una relazione di equivalenza su $a$ e sia $x \in a$.
    \begin{itemize}
        \item L'insieme $[x]_\sim := \{y \in a \mid x \sim y\}$ si chiama \textbf{classe di equivalenza di $x$ modulo $\sim$};
              ogni elemento di $[x]_\sim$ (ad esempio $x$ stesso) si dice \textbf{rappresentante} della classe.
        \item L'insieme di tutte le classi $a/{\sim} := \{[x]_\sim \mid x \in a\}$ si chiama \textbf{insieme quoziente} di $a$
              modulo $\sim$.
    \end{itemize}
    Quando non c'è ambiguità si scrive $[x]$ invece di $[x]_\sim$.
}

\begin{esempio}{Classi e quozienti}
    \begin{itemize}
        \item Per l'uguaglianza ogni classe è un singleton: $[x] = \{x\}$, e $a/{=} = \{\{x\} \mid x \in a\}$.
        \item Per la relazione totale su $a \neq \varnothing$ c'è una sola classe: $[x] = a$ per ogni $x$, e il quoziente è $\{a\}$.
        \item Per $\equiv_2$: $[0]_{\equiv_2}$ è l'insieme dei pari, $[1]_{\equiv_2}$ quello dei dispari, e $\mbZ/{\equiv_2}$ ha due elementi.
              Per $\equiv_3$ le classi sono i tre colori della figura nella spiegazione della
              Proposizione~\ref{prop:congruenza-equivalenza}: $[0] = \{\dots, -3, 0, 3, 6, \dots\}$, $[1] = \{\dots, -2, 1, 4, 7, \dots\}$,
              $[2] = \{\dots, -1, 2, 5, 8, \dots\}$.
        \item Per il nucleo di $n \in \mbZ \mapsto n^2 \in \mbZ$: $[n] = \{n, -n\}$; ad esempio $[0] = \{0\}$ e $[-3] = [3] = \{-3, 3\}$.
    \end{itemize}
\end{esempio}

\begin{concetto}
\begin{Teorema}{Proprietà fondamentali delle classi di equivalenza}{proprieta-classi}
    \begin{prerequisiti}
        \dip{Definizione}{def:classe-equivalenza}
        \dip{Definizione}{def:tipi-relazioni}
        \dip{Proposizione}{prop:doppia-inclusione}
        \dip{Definizione}{def:notazioni-assiomi}
        \dip{Proposizione}{prop:contrapposizione}
    \end{prerequisiti}
    Sia $\sim$ una relazione di equivalenza su $a$. Allora:
    \begin{enumerate}
        \item per ogni $x \in a$, $x \in [x]$; in particolare ogni classe è \textbf{non vuota};
        \item per ogni $x, y \in a$:\quad
              $x \sim y \ \leftrightarrow\ y \sim x \ \leftrightarrow\ y \in [x] \ \leftrightarrow\ x \in [y] \ \leftrightarrow\ [x] = [y]
              \ \leftrightarrow\ [x] \cap [y] \neq \varnothing$;
        \item per ogni $x, y \in a$, se $[x] \neq [y]$ allora $[x] \cap [y] = \varnothing$: le classi sono \textbf{a due a due disgiunte};
        \item $\bigcup (a/{\sim}) = a$.
    \end{enumerate}
\end{Teorema}
\begin{dimostrazione}
    \begin{caso}{Punto 1}
        Per la riflessività $x \sim x$, cioè $x \in [x]$.
    \end{caso}
    \begin{caso}{Punto 2}
        Le prime quattro condizioni sono equivalenti per la simmetria e per la definizione di classe
        ($y \in [x] \leftrightarrow x \sim y$, $x \in [y] \leftrightarrow y \sim x$). Chiudiamo il cerchio con le altre due.
        \begin{itemize}
            \item[] $x \sim y \rightarrow [x] = [y]$: se $w \in [x]$, cioè $x \sim w$, da $y \sim x$ (simmetria) e transitività si ha
                    $y \sim w$, cioè $w \in [y]$; quindi $[x] \subseteq [y]$. Scambiando $x$ e $y$, $[y] \subseteq [x]$: per la doppia
                    inclusione $[x] = [y]$.
            \item[] $[x] = [y] \rightarrow [x] \cap [y] \neq \varnothing$: per il punto 1 $x \in [x] = [y]$, quindi $x \in [x] \cap [y]$.
            \item[] $[x] \cap [y] \neq \varnothing \rightarrow x \sim y$: se $z \in [x] \cap [y]$, allora $x \sim z$ e $y \sim z$; per
                    simmetria $z \sim y$ e per transitività $x \sim y$.
        \end{itemize}
    \end{caso}
    \begin{caso}{Punto 3}
        Dal punto 2, $[x] \cap [y] \neq \varnothing \rightarrow [x] = [y]$; per contrapposizione $[x] \neq [y] \rightarrow [x] \cap [y] = \varnothing$.
    \end{caso}
    \begin{caso}{Punto 4}
        $(\subseteq)$ Gli elementi delle classi sono elementi di $a$. $(\supseteq)$ Se $x \in a$, per il punto 1 $x \in [x] \in a/{\sim}$,
        quindi $x \in \bigcup(a/{\sim})$.
    \end{caso}
\end{dimostrazione}
\end{concetto}

\begin{spiegazione}
    Le classi sono come le \textbf{stanze di una casa}: ogni elemento sta in una stanza (punto 1), due stanze diverse non hanno
    nessuno in comune (punto 3) e tutte insieme formano la casa intera (punto 4). Il punto 2 dice che \textbf{essere in
    relazione equivale a stare nella stessa stanza}, e che una stanza non dipende da chi scegliamo per indicarla:
    \begin{center}
    \begin{tikzpicture}
        \node[elem] (a1) at (0,1.1) {$1$}; \node[elem] (a2) at (1.0,1.35) {$2$}; \node[elem] (a3) at (0.55,0.35) {$3$};
        \node[elem] (b1) at (2.5,1.25) {$4$}; \node[elem] (b2) at (3.0,0.4) {$5$};
        \node[elem] (c1) at (4.4,0.85) {$6$};
        \blocco[7mm]{classeA}{(a1.center) -- (a2.center) -- (a3.center) -- cycle}
        \blocco[7mm]{classeB}{(b1.center) -- (b2.center)}
        \blocco[7mm]{classeC}{(c1.center) -- ++(0.001,0)}
        \draw[black!45, rounded corners=12pt] (-0.75,-0.45) rectangle (5.15,2.0);
        \node[font=\small] at (-0.45,1.75) {$a$};
        \node[font=\small, align=left, anchor=west] at (5.6,0.8)
            {$[1] = [2] = [3] = \{1, 2, 3\}$\\
             $[4] = [5] = \{4, 5\}$, \quad $[6] = \{6\}$\\
             $a/{\sim} = \{\{1, 2, 3\}, \{4, 5\}, \{6\}\}$};
    \end{tikzpicture}
    \end{center}
    Per dimostrare un'equivalenza tra tante condizioni non serve provare tutte le coppie di implicazioni: basta una catena che
    torni al punto di partenza, come abbiamo fatto con
    $x \sim y \Rightarrow [x] = [y] \Rightarrow [x] \cap [y] \neq \varnothing \Rightarrow x \sim y$.
\end{spiegazione}

Per il nucleo di equivalenza di una funzione le classi si descrivono con l'antiimmagine.

\begin{concetto}
\begin{Proposizione}{Classi del nucleo di equivalenza}{classi-nucleo}
    \begin{prerequisiti}
        \dip{Definizione}{def:nucleo}
        \dip{Definizione}{def:classe-equivalenza}
        \dip{Definizione}{def:immagine-antiimmagine}
        \dip{Assioma}{ass:estensionalita}
    \end{prerequisiti}
    Sia $f \colon a \to b$ e sia $x \in a$. Allora $[x]_{\sim_f} = \overleftarrow{f}(\{f(x)\})$.
\end{Proposizione}
\begin{dimostrazione}
    Per ogni $y \in a$:\quad $y \in [x]_{\sim_f} \leftrightarrow x \sim_f y \leftrightarrow f(y) = f(x) \leftrightarrow f(y) \in \{f(x)\}
    \leftrightarrow y \in \overleftarrow{f}(\{f(x)\})$. Per Estensionalità i due insiemi coincidono.
\end{dimostrazione}
\end{concetto}

\begin{spiegazione}
    La classe di $x$ è l'insieme di \textbf{tutti gli elementi che finiscono dove finisce $x$}: si parte da $f(x)$ e si
    seguono le frecce all'indietro. Nell'esempio di $f(n) = n^2$: $[2] = \overleftarrow{f}(\{4\}) = \{-2, 2\}$.
\end{spiegazione}

\dfn[proiezione-canonica]{Proiezione canonica}{
    Sia $\sim$ una relazione di equivalenza su $a$. La funzione
    \[ \pi \colon x \in a \mapsto [x]_\sim \in a/{\sim} \]
    si chiama \textbf{proiezione canonica} di $a$ su $a/{\sim}$.
}

\begin{concetto}
\begin{Proposizione}{Proprietà della proiezione canonica}{proiezione}
    \begin{prerequisiti}
        \dip{Definizione}{def:proiezione-canonica}
        \dip{Definizione}{def:suriettiva}
        \dip{Teorema}{teo:proprieta-classi}
        \dip{Definizione}{def:nucleo}
    \end{prerequisiti}
    Sia $\sim$ una relazione di equivalenza su $a$ e sia $\pi$ la proiezione canonica. Allora $\pi$ è suriettiva e il suo nucleo
    di equivalenza è proprio $\sim$: $\pi(x) = \pi(y) \leftrightarrow x \sim y$.
\end{Proposizione}
\begin{dimostrazione}
    \begin{caso}{Suriettiva}
        Ogni elemento di $a/{\sim}$ è una classe $[x]$ con $x \in a$, e $[x] = \pi(x)$.
    \end{caso}
    \begin{caso}{Il nucleo è $\sim$}
        $\pi(x) = \pi(y) \leftrightarrow [x] = [y] \leftrightarrow x \sim y$ per il Teorema~\ref{teo:proprieta-classi}, punto 2.
    \end{caso}
\end{dimostrazione}
\end{concetto}

\begin{spiegazione}
    $\pi$ «schiaccia» ogni stanza in un unico punto: manda ogni elemento nella sua classe. Nell'esempio delle stanze,
    $\pi(1) = \pi(2) = \pi(3) = \{1, 2, 3\}$. Il fatto che il nucleo di $\pi$ sia $\sim$ dice che \textbf{ogni relazione di
    equivalenza è il nucleo di qualche funzione}: non c'è differenza tra «dare un'equivalenza» e «dare una funzione e guardare
    quali elementi hanno la stessa immagine».
\end{spiegazione}

% ──────────────────────────────────────────────────────────────
\section{Il teorema fondamentale di omomorfismo per insiemi}\label{sec:omomorfismo-insiemi}

Ogni funzione si può spezzare in tre passi elementari. È il primo esempio di un'idea che ritroveremo spesso: quella di
\textbf{passare al quoziente}.

\begin{concetto}
\begin{Teorema}{Teorema fondamentale di omomorfismo per insiemi}{omomorfismo-insiemi}
    \begin{prerequisiti}
        \dip{Definizione}{def:nucleo}
        \dip{Definizione}{def:proiezione-canonica}
        \dip{Definizione}{def:funzioni-notevoli}
        \dip{Definizione}{def:funzione}
        \dip{Definizione}{def:biettiva}
        \dip{Teorema}{teo:proprieta-classi}
        \dip{Corollario}{cor:terna}
    \end{prerequisiti}
    Sia $f \colon a \to b$ una funzione, siano $\pi \colon a \to a/{\sim_f}$ la proiezione canonica e $\iota \colon \mathrm{Im} f \to b$
    l'immersione. Allora:
    \begin{enumerate}
        \item la corrispondenza $\tilde{f} \colon [x]_{\sim_f} \in a/{\sim_f} \mapsto f(x) \in \mathrm{Im} f$ è una funzione \textbf{biettiva};
        \item $f = \iota \circ \tilde{f} \circ \pi$, cioè il diagramma seguente è commutativo:
    \end{enumerate}
    \begin{center}
    \begin{tikzcd}[column sep=2.4cm, row sep=1.1cm]
        a \arrow[r, "f"] \arrow[d, "\pi"'] & b \\
        a/{\sim_f} \arrow[r, "\tilde{f}"'] & \mathrm{Im} f \arrow[u, hook, "\iota"']
    \end{tikzcd}
    \end{center}
\end{Teorema}
\begin{dimostrazione}
    \begin{caso}{$\tilde{f}$ è una funzione (è «ben posta»)}
        $\tilde{f}$ è la corrispondenza $(a/{\sim_f}, \mathrm{Im} f, \tilde{g})$ con $\tilde{g} = \{([x], f(x)) \mid x \in a\}$. Ogni classe
        ha un'immagine, perché ogni classe è $[x]$ per qualche $x$. L'immagine è unica: se $[x] = [y]$, allora $x \sim_f y$, cioè
        $f(x) = f(y)$. Quindi \textbf{il valore non dipende dal rappresentante} scelto.
    \end{caso}
    \begin{caso}{$\tilde{f}$ è iniettiva}
        Se $\tilde{f}([x]) = \tilde{f}([y])$, allora $f(x) = f(y)$, cioè $x \sim_f y$, e quindi $[x] = [y]$
        (Teorema~\ref{teo:proprieta-classi}, punto 2).
    \end{caso}
    \begin{caso}{$\tilde{f}$ è suriettiva}
        Se $w \in \mathrm{Im} f$, esiste $x \in a$ con $f(x) = w$, e allora $\tilde{f}([x]) = w$.
    \end{caso}
    \begin{caso}{$f = \iota \circ \tilde{f} \circ \pi$}
        Le due funzioni hanno lo stesso dominio $a$ e lo stesso codominio $b$, e per ogni $x \in a$:
        \[ \iota\big(\tilde{f}(\pi(x))\big) = \iota\big(\tilde{f}([x])\big) = \iota(f(x)) = f(x). \]
        Hanno quindi lo stesso grafico e sono la stessa terna (Corollario~\ref{cor:terna}).
    \end{caso}
\end{dimostrazione}
\end{concetto}

\begin{spiegazione}
    Prendiamo $a = \{1, 2, 3, 4, 5\}$, $b = \{p, q, r, s\}$ e $f$ con $f(1) = f(2) = p$, $f(3) = q$, $f(4) = f(5) = r$. Allora $f$
    fa tre cose, una dopo l'altra:
    \begin{center}
    \begin{tikzpicture}
        \foreach \e/\y in {1/2.0, 2/1.5, 3/0.5, 4/-0.5, 5/-1.0} \node[elem] (x\e) at (0,\y) {$\e$};
        \blocco[6mm]{classeA}{(x1.center) -- (x2.center)}
        \blocco[6mm]{classeB}{(x3.center) -- ++(0.001,0)}
        \blocco[6mm]{classeC}{(x4.center) -- (x5.center)}
        \node[elem] (c1) at (3.1,1.75) {$\{1, 2\}$};
        \node[elem] (c2) at (3.1,0.5) {$\{3\}$};
        \node[elem] (c3) at (3.1,-0.75) {$\{4, 5\}$};
        \node[elem] (i1) at (6.0,1.75) {$p$}; \node[elem] (i2) at (6.0,0.5) {$q$}; \node[elem] (i3) at (6.0,-0.75) {$r$};
        \node[elem] (b1) at (8.6,1.75) {$p$}; \node[elem] (b2) at (8.6,0.5) {$q$}; \node[elem] (b3) at (8.6,-0.75) {$r$};
        \node[elem] (b4) at (8.6,-1.7) {$s$};
        \draw[fun3] (x1) -- (c1); \draw[fun3] (x2) -- (c1); \draw[fun3] (x3) -- (c2); \draw[fun3] (x4) -- (c3); \draw[fun3] (x5) -- (c3);
        \draw[fun] (c1) -- (i1); \draw[fun] (c2) -- (i2); \draw[fun] (c3) -- (i3);
        \draw[fun2, {Hooks[right]}-{Stealth[length=2.2mm]}] (i1) -- (b1);
        \draw[fun2, {Hooks[right]}-{Stealth[length=2.2mm]}] (i2) -- (b2);
        \draw[fun2, {Hooks[right]}-{Stealth[length=2.2mm]}] (i3) -- (b3);
        \node[font=\small] at (0,2.75) {$a$};
        \node[font=\small] at (3.1,2.75) {$a/{\sim_f}$};
        \node[font=\small] at (6.0,2.75) {$\mathrm{Im} f$};
        \node[font=\small] at (8.6,2.75) {$b$};
        \node[font=\small, seagreen!60!black] at (1.55,2.75) {$\pi$};
        \node[font=\small, primary!80!black] at (4.55,2.75) {$\tilde{f}$};
        \node[font=\small, cherryred!85!black] at (7.3,2.75) {$\iota$};
    \end{tikzpicture}
    \end{center}
    \begin{passi}
        \item $\pi$ \textbf{raggruppa} gli elementi che hanno la stessa immagine: $1$ e $2$ finiscono nella stessa classe.
              È suriettiva, ma in generale non iniettiva.
        \item $\tilde{f}$ dà a ogni gruppo il suo valore: ora gruppi diversi hanno valori diversi, quindi è \textbf{iniettiva}; e
              arriva solo sui valori effettivamente presi ($\mathrm{Im} f$), quindi è anche \textbf{suriettiva}.
        \item $\iota$ \textbf{rimette} i valori nel codominio $b$, dove c'è anche $s$, che nessuno raggiunge. È iniettiva, ma
              in generale non suriettiva.
    \end{passi}
    Tutto ciò che impedisce a $f$ di essere iniettiva è finito in $\pi$, tutto ciò che le impedisce di essere suriettiva è
    finito in $\iota$; in mezzo resta una biezione.
\end{spiegazione}

\begin{concetto}
\begin{Corollario}{Ogni funzione è composta di una suriettiva e di un'iniettiva}{iniettiva-suriettiva}
    \begin{prerequisiti}
        \dip{Teorema}{teo:omomorfismo-insiemi}
        \dip{Proposizione}{prop:proiezione}
        \dip{Teorema}{teo:comp-iniettive}
    \end{prerequisiti}
    Ogni funzione $f \colon a \to b$ si scrive come $f = j \circ \pi$, con $\pi$ suriettiva e $j$ iniettiva.
\end{Corollario}
\begin{dimostrazione}
    Per il Teorema~\ref{teo:omomorfismo-insiemi}, $f = \iota \circ \tilde{f} \circ \pi = (\iota \circ \tilde{f}) \circ \pi$. La proiezione
    $\pi$ è suriettiva (Proposizione~\ref{prop:proiezione}); $j = \iota \circ \tilde{f}$ è iniettiva perché composta di due funzioni
    iniettive (Teorema~\ref{teo:comp-iniettive}): $\tilde{f}$ è biettiva e l'immersione $\iota$ è iniettiva.
\end{dimostrazione}
\end{concetto}

\begin{spiegazione}
    Con $f \colon n \in \mbZ \mapsto n^2 \in \mbZ$:
    \begin{passi}
        \item $\pi \colon n \mapsto [n] = \{n, -n\}$ raggruppa ogni intero con il suo opposto (suriettiva su
              $\mbZ/{\sim_f} = \{\{0\}, \{1, -1\}, \{2, -2\}, \dots\}$);
        \item $j \colon \{n, -n\} \mapsto n^2 \in \mbZ$ manda ogni coppia nel suo quadrato: coppie diverse hanno quadrati diversi,
              quindi è iniettiva (ma non suriettiva: $2$ non è un quadrato).
    \end{passi}
    Infatti $j(\pi(-3)) = j(\{3, -3\}) = 9 = f(-3)$.
\end{spiegazione}

% ──────────────────────────────────────────────────────────────
\section{Partizioni}\label{sec:partizioni}

Le proprietà fondamentali delle classi (Teorema~\ref{teo:proprieta-classi}) dicono che un insieme quoziente divide l'insieme
in pezzi non vuoti, disgiunti, che lo ricoprono tutto. Diamo un nome a questo tipo di suddivisione.

\dfn[partizione]{Partizione}{
    Sia $a$ un insieme. Un insieme $p$ di parti di $a$ (cioè $p \subseteq \mcP(a)$) si dice \textbf{partizione} di $a$ se
    \begin{enumerate}
        \item $\forall x \in p\,(x \neq \varnothing)$: i pezzi sono non vuoti;
        \item $\forall x, y \in p\,(x \neq y \rightarrow x \cap y = \varnothing)$: i pezzi sono a due a due disgiunti;
        \item $\bigcup p = a$: i pezzi ricoprono $a$.
    \end{enumerate}
    Gli elementi di $p$ si chiamano anche \textbf{blocchi} della partizione. Indichiamo con $\mathrm{Part}(a)$ l'insieme delle
    partizioni di $a$ (è una parte di $\mcP(\mcP(a))$).
}

\begin{esempio}{Le cinque partizioni di $\{0, 1, 2\}$}
    \begin{center}
    \begin{tikzpicture}
        \foreach \k/\px in {1/0, 2/3.0, 3/6.0, 4/9.0, 5/12.0} {
            \begin{scope}[xshift=\px cm]
                \node[elem] (z\k) at (0.55,0.95) {$0$};
                \node[elem] (u\k) at (0,0) {$1$};
                \node[elem] (d\k) at (1.1,0) {$2$};
            \end{scope}
        }
        \blocco{classeA}{(z1.center) -- ++(0.001,0)} \blocco{classeB}{(u1.center) -- ++(0.001,0)} \blocco{classeC}{(d1.center) -- ++(0.001,0)}
        \blocco{classeA}{(z2.center) -- (u2.center)} \blocco{classeC}{(d2.center) -- ++(0.001,0)}
        \blocco{classeA}{(z3.center) -- (d3.center)} \blocco{classeB}{(u3.center) -- ++(0.001,0)}
        \blocco{classeB}{(u4.center) -- (d4.center)} \blocco{classeA}{(z4.center) -- ++(0.001,0)}
        \blocco{classeD}{(z5.center) -- (u5.center) -- (d5.center) -- cycle}
        \node[font=\scriptsize, align=center] at (0.55,-0.85) {$\{\{0\}, \{1\}, \{2\}\}$};
        \node[font=\scriptsize, align=center] at (3.55,-0.85) {$\{\{0, 1\}, \{2\}\}$};
        \node[font=\scriptsize, align=center] at (6.55,-0.85) {$\{\{0, 2\}, \{1\}\}$};
        \node[font=\scriptsize, align=center] at (9.55,-0.85) {$\{\{1, 2\}, \{0\}\}$};
        \node[font=\scriptsize, align=center] at (12.55,-0.85) {$\{\{0, 1, 2\}\}$};
    \end{tikzpicture}
    \end{center}
    Invece $\{\{0, 1\}, \{1, 2\}\}$ non è una partizione (i blocchi non sono disgiunti), $\{\{0\}, \{1\}\}$ non lo è (non ricopre $2$)
    e $\{\{0, 1, 2\}, \varnothing\}$ nemmeno (contiene un blocco vuoto). Altri esempi: $\{\text{pari}, \text{dispari}\}$ è una
    partizione di $\mbZ$; l'unica partizione di $\varnothing$ è $\varnothing$ (le tre condizioni sono vere a vuoto).
\end{esempio}

\begin{concetto}
\begin{Corollario}{Ogni insieme quoziente è una partizione}{quoziente-partizione}
    \begin{prerequisiti}
        \dip{Teorema}{teo:proprieta-classi}
        \dip{Definizione}{def:partizione}
    \end{prerequisiti}
    Se $\sim$ è una relazione di equivalenza su $a$, allora $a/{\sim} \in \mathrm{Part}(a)$.
\end{Corollario}
\begin{dimostrazione}
    Ogni classe è una parte di $a$, quindi $a/{\sim} \subseteq \mcP(a)$. Le tre condizioni della Definizione~\ref{def:partizione}
    sono i punti 1, 3 e 4 del Teorema~\ref{teo:proprieta-classi}.
\end{dimostrazione}
\end{concetto}

\begin{spiegazione}
    È la casa con le stanze vista dall'alto: le stanze (le classi) sono i blocchi della partizione. Ad esempio $\mbZ/{\equiv_3}$
    è la partizione di $\mbZ$ nei tre colori della figura sulla congruenza modulo $3$.
\end{spiegazione}

Vale anche il viceversa: ogni partizione è l'insieme quoziente di una (e una sola) relazione di equivalenza. Le relazioni di
equivalenza e le partizioni sono quindi \textbf{due modi di dare la stessa informazione}.

\begin{concetto}
\begin{Teorema}{Teorema fondamentale su relazioni di equivalenza e partizioni}{equivalenze-partizioni}
    \begin{prerequisiti}
        \dip{Corollario}{cor:quoziente-partizione}
        \dip{Teorema}{teo:proprieta-classi}
        \dip{Definizione}{def:partizione}
        \dip{Definizione}{def:biettiva}
        \dip{Corollario}{cor:terna}
    \end{prerequisiti}
    Per ogni insieme $a$ la funzione
    \[ \Phi \colon {\sim} \in \mathrm{Eq}(a) \mapsto a/{\sim} \in \mathrm{Part}(a) \]
    è biettiva. La sua inversa associa alla partizione $p$ la relazione «stare nello stesso blocco»:
    $x \sim_p y \leftrightarrow \exists w \in p\,(x \in w \wedge y \in w)$.
\end{Teorema}
\begin{dimostrazione}
    $\Phi$ è una funzione a valori in $\mathrm{Part}(a)$ per il Corollario~\ref{cor:quoziente-partizione}.
    \begin{caso}{$\Phi$ è iniettiva}
        Siano $\sim_1, \sim_2 \in \mathrm{Eq}(a)$ con $a/{\sim_1} = a/{\sim_2}$, e sia $x \sim_1 y$. Allora $[x]_{\sim_1} = [y]_{\sim_1}$
        è un elemento di $a/{\sim_1} = a/{\sim_2}$, quindi esiste $z \in a$ con $[x]_{\sim_1} = [z]_{\sim_2}$. Perciò $x, y \in [z]_{\sim_2}$,
        cioè $z \sim_2 x$ e $z \sim_2 y$, e per simmetria e transitività $x \sim_2 y$. Scambiando i ruoli, $x \sim_2 y \rightarrow x \sim_1 y$.
        Le due relazioni hanno lo stesso grafico, quindi sono uguali (Corollario~\ref{cor:terna}).
    \end{caso}
    \begin{caso}{$\Phi$ è suriettiva}
        Sia $p \in \mathrm{Part}(a)$ e sia $\sim_p$ come nell'enunciato.
        \begin{itemize}
            \item[] \textbf{Riflessiva}: poiché $\bigcup p = a$, ogni $x \in a$ sta in un blocco $w$, quindi $x \sim_p x$.
            \item[] \textbf{Simmetrica}: $x \in w \wedge y \in w$ equivale a $y \in w \wedge x \in w$.
            \item[] \textbf{Transitiva}: se $x, y \in w_1$ e $y, z \in w_2$ con $w_1, w_2 \in p$, allora $y \in w_1 \cap w_2$; i blocchi
                    sono a due a due disgiunti, quindi $w_1 = w_2$ e $x, z$ stanno nello stesso blocco.
        \end{itemize}
        Resta da vedere che $a/{\sim_p} = p$. Sia $x \in a$ e sia $w_x$ il blocco che contiene $x$ (è unico, perché i blocchi sono
        disgiunti). Allora $y \in [x]_{\sim_p} \leftrightarrow$ $x$ e $y$ stanno in uno stesso blocco $\leftrightarrow y \in w_x$, cioè
        $[x]_{\sim_p} = w_x$. Quindi ogni classe è un blocco; e ogni blocco $w$, non essendo vuoto, contiene un $x$ ed è $w = w_x = [x]_{\sim_p}$.
        Dunque $\Phi(\sim_p) = a/{\sim_p} = p$.
    \end{caso}
\end{dimostrazione}
\end{concetto}

\begin{spiegazione}
    Ecco la stessa informazione nelle due forme, su $a = \{1, 2, 3, 4, 5\}$. A sinistra il grafo della relazione di
    equivalenza (le doppie frecce indicano andata e ritorno; ogni elemento ha anche il suo cappio, che omettiamo); a destra
    la partizione corrispondente:
    \begin{center}
    \begin{tikzpicture}
        \node[rv] (1) at (0,1.2) {$1$}; \node[rv] (2) at (1.3,1.2) {$2$}; \node[rv] (3) at (0.65,0) {$3$};
        \node[rv] (4) at (2.7,1.2) {$4$}; \node[rv] (5) at (2.7,0) {$5$};
        \draw[{Stealth[length=2mm]}-{Stealth[length=2mm]}, line width=0.7pt, black!70] (1) -- (2);
        \draw[{Stealth[length=2mm]}-{Stealth[length=2mm]}, line width=0.7pt, black!70] (2) -- (3);
        \draw[{Stealth[length=2mm]}-{Stealth[length=2mm]}, line width=0.7pt, black!70] (1) -- (3);
        \draw[{Stealth[length=2mm]}-{Stealth[length=2mm]}, line width=0.7pt, black!70] (4) -- (5);
        \node[font=\small] at (1.35,-0.8) {la relazione $\sim$};
        \draw[-{Stealth}, line width=1pt, primary] (3.6,0.6) -- node[above, font=\small] {$\Phi$} (5.0,0.6);
        \draw[{Stealth}-, line width=1pt, seagreen!65!black] (3.6,0.2) -- node[below, font=\small] {$\Phi^{-1}$} (5.0,0.2);
        \begin{scope}[xshift=5.9cm]
            \node[elem] (p1) at (0,1.2) {$1$}; \node[elem] (p2) at (1.3,1.2) {$2$}; \node[elem] (p3) at (0.65,0) {$3$};
            \node[elem] (p4) at (2.7,1.2) {$4$}; \node[elem] (p5) at (2.7,0) {$5$};
            \blocco[7mm]{classeA}{(p1.center) -- (p2.center) -- (p3.center) -- cycle}
            \blocco[7mm]{classeB}{(p4.center) -- (p5.center)}
            \node[font=\small] at (1.35,-0.8) {la partizione $\{\{1, 2, 3\}, \{4, 5\}\}$};
        \end{scope}
    \end{tikzpicture}
    \end{center}
    \begin{passi}
        \item Da sinistra a destra ($\Phi$): ogni «gruppo di elementi tutti collegati tra loro» diventa un blocco.
        \item Da destra a sinistra ($\Phi^{-1}$): si collegano tra loro (con cappi e frecce in entrambi i versi) gli elementi di ogni
              blocco, e solo quelli.
        \item L'iniettività dice che due relazioni diverse danno partizioni diverse; la suriettività che ogni partizione si
              ottiene così.
    \end{passi}
\end{spiegazione}

% ──────────────────────────────────────────────────────────────
\section{Trovare tutte le relazioni di equivalenza}\label{sec:trovare-equivalenze}

Grazie al Teorema~\ref{teo:equivalenze-partizioni}, \textbf{trovare tutte le relazioni di equivalenza su un insieme equivale a
trovare tutte le sue partizioni}, che è molto più facile: basta disegnare i blocchi.

\begin{esempio}{Le relazioni di equivalenza su $\{0, 1, 2\}$}
    Le partizioni di $\{0, 1, 2\}$ sono le cinque disegnate prima, quindi le relazioni di equivalenza sono cinque. I loro grafici:
    \begin{center}
    \renewcommand{\arraystretch}{1.3}
    \begin{tabular}{l|l}
        \rowcolor{primary!12} \textbf{Partizione} & \textbf{Grafico della relazione} (con $D = \mathrm{diag}(\{0, 1, 2\})$) \\
        $\{\{0\}, \{1\}, \{2\}\}$ & $D$ \quad (l'uguaglianza) \\
        $\{\{0, 1\}, \{2\}\}$ & $D \cup \{(0, 1), (1, 0)\}$ \\
        $\{\{0, 2\}, \{1\}\}$ & $D \cup \{(0, 2), (2, 0)\}$ \\
        $\{\{1, 2\}, \{0\}\}$ & $D \cup \{(1, 2), (2, 1)\}$ \\
        $\{\{0, 1, 2\}\}$ & $\{0, 1, 2\} \times \{0, 1, 2\}$ \quad (la relazione totale) \\
    \end{tabular}
    \end{center}
\end{esempio}

\begin{esempio}{Equivalenze con dei vincoli}
    Siano $a, b, c, d$ elementi a due a due distinti e $s = \{a, b, c, d\}$.
    \begin{itemize}
        \item Le equivalenze su $s$ con $[a] = [b] = [c]$ corrispondono alle partizioni in cui $a, b, c$ stanno nello stesso blocco:
              $\{\{a, b, c\}, \{d\}\}$ e $\{\{a, b, c, d\}\}$. Sono \textbf{due}.
        \item Quelle con $[a] = [b]$ corrispondono alle partizioni con $a$ e $b$ nello stesso blocco: si parte da $\{a, b\}$ e si
              decide dove mettere $c$ e $d$. Si ottengono $\{\{a, b\}, \{c\}, \{d\}\}$, $\{\{a, b\}, \{c, d\}\}$, $\{\{a, b, c\}, \{d\}\}$,
              $\{\{a, b, d\}, \{c\}\}$ e $\{\{a, b, c, d\}\}$: sono \textbf{cinque}.
    \end{itemize}
\end{esempio}

\begin{esempio}{Un esempio passo passo}
    Sia $a = \{n \in \mbN \mid n \leq 7\}$. Troviamo tutte le relazioni di equivalenza $\rho$ su $a$ tali che
    $0 \rho 7$, $1 \rho 4$, $\{3, 4, 5, 7\} \subseteq [2]_\rho$ e $0 \rho 6$. Ogni vincolo costringe certi elementi nello stesso blocco:
    \begin{passi}
        \item $0 \rho 7$: il blocco di $0$ contiene $\{0, 7\}$.
        \item $1 \rho 4$: il blocco di $1$ contiene $\{1, 4\}$.
        \item $\{3, 4, 5, 7\} \subseteq [2]$: $2, 3, 4, 5, 7$ stanno nello stesso blocco. Ma $7$ è con $0$ e $4$ è con $1$, quindi
              lo stesso blocco contiene $\{0, 1, 2, 3, 4, 5, 7\}$.
        \item $0 \rho 6$: anche $6$ entra nel blocco, che diventa tutto $a$.
    \end{passi}
    C'è quindi \textbf{una sola} relazione con queste proprietà: la relazione totale su $a$.
\end{esempio}
